两个多项式相乘就是系数序列的普通卷积：
\[
  c_k=\sum_{i+j=k}a_i b_j.
\]
若两多项式长度分别为 $n,m$，结果长度为 $n+m-1$；FFT 和 NTT 都需要把变换长度补到不小于它的二次幂。

\subsubsection{FFT}

FFT 在复数域上计算普通卷积，代码使用迭代蝶形，逆变换后用 \texttt{llround} 还原整数系数。
复杂度为 $O(N\log N)$，其中 $N$ 是补齐后的长度。
它不依赖特定模数，但系数或长度过大时会产生浮点误差；此时应改用 NTT、多个 NTT 合并或系数拆分。

\subsubsection{NTT}

NTT 在有限域上计算普通卷积。
代码固定使用
\[
  P=998244353=119\cdot2^{23}+1,
  \qquad G=3,
\]
其中 $G$ 是原根，故变换长度必须是不超过 $2^{23}$ 的二次幂。
每层现场计算单位根，不额外维护根表或位逆序表，复杂度为 $O(N\log N)$。

\subsubsection{FWT}

FWT 处理下标之间的按位卷积：
\[
  c_s^{\mathrm{OR}}
  =\sum_{i\mathbin{|}j=s}a_i b_j,
  \qquad
  c_s^{\mathrm{AND}}
  =\sum_{i\mathbin{\&}j=s}a_i b_j,
\]
\[
  c_s^{\mathrm{XOR}}
  =\sum_{i\mathbin{\oplus}j=s}a_i b_j.
\]
它们不是子集卷积。
数组长度应为 $2^d$，不足时补零；若两数组长度不同，补到共同的二次幂长度。

设一个蝶形中的两个数为 $(x,y)$，正变换分别为
\[
\begin{aligned}
  \mathrm{OR}:  &(x,y)\mapsto(x,x+y),\\
  \mathrm{AND}: &(x,y)\mapsto(x+y,y),\\
  \mathrm{XOR}: &(x,y)\mapsto(x+y,x-y).
\end{aligned}
\]
逆变换分别为
\[
\begin{aligned}
  \mathrm{OR}:  &(x,y)\mapsto(x,y-x),\\
  \mathrm{AND}: &(x,y)\mapsto(x-y,y),\\
  \mathrm{XOR}: &(x,y)\mapsto
    \left(\frac{x+y}{2},\frac{x-y}{2}\right).
\end{aligned}
\]
代码在模 $998244353$ 下运算，因此 XOR 逆变换中的 $2$ 可逆。

若 $*$ 表示选定类型的按位卷积，则变换后卷积变成逐点乘法：
\[
  \widehat{A*B}_i=\widehat A_i\widehat B_i.
\]
因此 $A$ 的 $k$ 次卷积幂满足
\[
  \widehat{A^{*k}}_i=(\widehat A_i)^k.
\]
只需一次正变换、每个点做一次快速幂、再做一次逆变换，复杂度为
\[
  O(N\log N+N\log k).
\]
当 $k=0$ 时结果是卷积单位元：OR、XOR 的单位元位于下标 $0$，AND 的单位元位于下标 $N-1$。
